\section{Introduction}
\label{sec:introduction}

This survey examines regional English accent as an important factor in
speech-to-text translation. It focuses on two main system types: Direct speech
translation (Direct ST) and Cascaded speech translation (Cascaded ST). Direct
ST translates source-language speech directly into target-language text
without producing an intermediate transcript. Cascaded ST first uses automatic
speech recognition (ASR) to convert speech into text and then applies machine
translation (MT) to produce the target-language translation. These two
approaches are compared throughout the survey, without assuming that either
one is always more accurate or more robust.

The survey is guided by the following research question:
\begin{quote}
\emph{How does regional English accent affect the translation quality of
Direct and Cascaded speech-to-text translation systems?}
\end{quote}
The main focus is final translation quality across regional accent groups.
Word error rate (WER) is considered only as a diagnostic measure for the ASR
stage of the Cascaded system. It is not treated as a direct measure of
translation quality.

The comparison is important because Direct and Cascaded systems have different
possible failure points. Cascaded systems can benefit from mature ASR and MT
models and from the larger datasets available for these separate tasks.
However, errors in the ASR transcript can be passed to the MT stage. Direct
systems remove this intermediate transcript, but they depend more strongly on
paired speech--translation data \cite{sperber2020endtoend}. Existing
comparisons do not show that one system type is always better. Results vary
depending on language pair, training resources, test segmentation, reference
translations, and evaluation metrics
\cite{bentivogli2021cascade,anastasopoulos2021iwslt,
etchegoyhen2022casestudy}. Therefore, overall system performance does not
directly show whether one approach is more robust to a particular regional
accent.

Accent is an important evaluation factor because ASR performance is not equal
across English varieties. Previous studies have reported performance
differences across regional, dialectal, and second-language accents, including
poorer performance on accents that were not represented during training
\cite{dichristofano2022global,jahan2025unveiling,shi2021aesrc,
prabhu2023accentcodebooks}. However, these findings mainly show an ASR-level
problem rather than a final translation result. This distinction is important
because similar WER values can lead to different levels of translation error,
while Direct ST does not produce an intermediate transcript where WER can be
measured \cite{le2017disentangling}. Previous work also shows that phonetic
recognition errors can reduce translation quality in Cascaded systems
\cite{sperber2019,pida2026}. This makes an accent effect on translation
plausible, but such an effect must still be measured using final translations
rather than inferred from ASR scores alone.

The literature was selected by theme rather than organised as a paper-by-paper
catalogue. The search covered Direct and Cascaded ST, ASR-to-MT error
propagation, accented ASR, unseen-accent generalisation, accented or dialectal
speech translation, accent-labelled datasets, speech-translation datasets,
evaluation metrics, and statistical reliability. Candidate papers and
publication details were checked through official sources, including ACL
Anthology, ISCA Archive, and IEEE proceedings. ArXiv papers were included when
no published venue version was identified. A paper was retained when its full
text had been reviewed and it directly contributed to at least one theme,
system comparison, evaluation issue, or test of the proposed research gap.
Blogs, secondary summaries, duplicate sources, and papers without verifiable
support for the required claims were excluded.

This process retained 27 studies, including recent work and older foundational
papers on evaluation metrics. Because a common systematic search record was
not maintained across all contributors, the survey is treated as a focused
narrative literature survey rather than an exhaustive systematic review.

The survey makes three main contributions. First, it connects research on
Direct and Cascaded speech translation with research on accent-sensitive ASR,
while keeping ASR error and translation quality as separate concepts. Second,
it examines whether the proposed research gap has already been addressed by
the closest studies on accent robustness, speech translation, and error
propagation. Third, it identifies the dataset, reference, metric, and
statistical controls needed for a fair comparison of final translation quality
across regional accent groups.

The remainder of the survey is organised as follows. Section
\ref{sec:architectures} reviews Direct and Cascaded speech-translation systems
and their empirical comparisons. Section~\ref{sec:accent-asr} examines accent
robustness in ASR. Section~\ref{sec:accent-st} reviews work that connects
accent, error propagation, and speech translation. Section~\ref{sec:evaluation}
discusses datasets, reference quality, evaluation metrics, and statistical
reliability. Section~\ref{sec:synthesis} combines the findings across these
themes to assess the research gap. Finally, Section~\ref{sec:conclusion}
presents the conclusion and the research question supported by the literature.
